\documentclass[11pt,a4paper]{article}
\usepackage[utf8]{inputenc}
\usepackage[T1]{fontenc}
\IfFileExists{french.ldf}{\usepackage[french]{babel}}{\usepackage[english]{babel}}
\usepackage{amsmath,amssymb,amsthm}
\usepackage{geometry}\geometry{margin=2.3cm}
\usepackage{listings}
\usepackage{xcolor}
\usepackage{enumitem}
\usepackage{fancyhdr}
\usepackage{lmodern}
\usepackage{titlesec}
\usepackage{tikz}
\usetikzlibrary{arrows.meta,positioning}
\usepackage{hyperref}
\hypersetup{colorlinks=true, urlcolor=[RGB]{2,101,115}, linkcolor=black}

\definecolor{codebg}{RGB}{247,249,251}
\definecolor{kw}{RGB}{175,45,45}
\definecolor{cmt}{RGB}{110,120,130}
\definecolor{str}{RGB}{20,110,60}
\definecolor{cnc}{RGB}{2,101,115}
\lstset{
  language=Python, backgroundcolor=\color{codebg}, basicstyle=\ttfamily\small,
  keywordstyle=\color{kw}\bfseries, commentstyle=\color{cmt}\itshape, stringstyle=\color{str},
  numbers=left, numberstyle=\tiny\color{cmt}, numbersep=8pt, frame=single, rulecolor=\color{gray!40},
  showstringspaces=false, breaklines=true, columns=fullflexible, extendedchars=true, inputencoding=utf8,
  literate={é}{{\'e}}1 {è}{{\`e}}1 {à}{{\`a}}1 {ê}{{\^e}}1 {î}{{\^i}}1 {ç}{{\c c}}1 {ù}{{\`u}}1 {ô}{{\^o}}1 {û}{{\^u}}1 {'}{{'}}1
}

\newcounter{qcnt}
\newcommand{\Q}{\medskip\par\stepcounter{qcnt}\noindent\textbf{Question~\theqcnt.}~}
\newcommand{\code}[1]{\lstinline!#1!}
\newcommand{\hh}[1]{h(#1)}

\newtheorem{lemme}{Lemme}
\newtheorem{theoreme}{Théorème}
\theoremstyle{definition}
\newtheorem{defi}{Définition}

\pagestyle{fancy}\fancyhf{}
\lhead{\small Préparation au CNC --- Informatique MP}
\rhead{\small Protocoles TCP/IP et HTTP/HTTPS}
\cfoot{\thepage}
\renewcommand{\headrulewidth}{0.4pt}
\titleformat*{\section}{\large\bfseries\scshape\color{cnc}}
\titleformat*{\subsection}{\bfseries\color{cnc}}
\usepackage{longtable}
\newcommand{\rfc}[1]{\href{https://www.rfc-editor.org/rfc/rfc#1}{\textsc{rfc}~#1}}

\begin{document}
\begin{center}
{\large\bfseries Préparation au Concours National Commun --- Informatique MP}\\[8pt]
\rule{0.9\linewidth}{0.5pt}\\[6pt]
{\Large\bfseries Épreuve d'Informatique}\\[4pt]
{\bfseries Filière MP --- Durée : 4 heures}\\[6pt]
{\itshape Voyage d'un paquet : de l'adresse IP au cadenas du navigateur}\\[4pt]
{\itshape Les protocoles TCP/IP et HTTP/HTTPS}\\[6pt]
\rule{0.9\linewidth}{0.5pt}\\[8pt]
{\small Auteur~: \href{https://orcid.org/0000-0002-6108-7176}{\textbf{Pr Agrégé El Hadiq Zouhair}}}\\[3pt]
{\small \href{https://cpgeacademy.org}{\textbf{cpgeacademy.org}} \quad---\quad Tous droits réservés \textcopyright\ cpgeacademy.org}
\end{center}
\medskip
\begin{center}
\fcolorbox{gray}{gray!5}{\begin{minipage}{0.93\linewidth}\small
\textbf{Avis important.} Ce document est une \textbf{initiative personnelle} du professeur,
à but strictement pédagogique. Il ne s'agit \textbf{pas} d'un sujet officiel du Concours
National Commun~; il s'en inspire uniquement par la forme et l'esprit.
\smallskip

\textbf{Sources.} Le sujet s'appuie sur les documents normatifs de l'\textsc{ietf} et sur
deux articles fondateurs~: \rfc{791} (IP), \rfc{4632} (CIDR), \rfc{1071} et \rfc{1624}
(somme de contrôle), \rfc{6298} (temporisateur de retransmission), \rfc{5681} (contrôle de
congestion), \rfc{9110} et \rfc{9112} (HTTP/1.1), \rfc{9111} (cache), \rfc{8446}
(TLS~1.3)~; V.~Jacobson, \emph{Congestion Avoidance and Control} (SIGCOMM~1988) et
M.~Mathis \emph{et al.}, \emph{The Macroscopic Behavior of the TCP Congestion Avoidance
Algorithm} (CCR~1997). \textbf{Aucune connaissance préalable en réseaux n'est requise}~:
tout ce qui est nécessaire est défini dans le sujet.
\smallskip

\textbf{Pour aller plus loin.} Les élèves peuvent traiter une \textbf{nouvelle version} de
ce problème dotée d'un \textbf{autocorrecteur des réponses}, et suivre un \textbf{cours
complet}, sur \href{https://cpgeacademy.org}{\textbf{cpgeacademy.org}}.
\end{minipage}}
\end{center}
\medskip


\section*{Présentation générale}

Lorsqu'un navigateur affiche une page, une longue chaîne de mécanismes s'est mise en marche~:
l'adresse de destination a été comparée à une table de routage, chaque paquet a été muni
d'une somme de contrôle, un protocole de transport a garanti que rien n'a été perdu ni
dupliqué, un texte a été analysé syntaxiquement, un cache a évité de retélécharger ce qui
était déjà connu, et un échange cryptographique a établi un secret partagé.

Ce problème étudie \textbf{algorithmiquement} chacun de ces mécanismes. Il ne suppose
\emph{aucune} connaissance en réseaux~: tous les objets manipulés sont définis dans le
sujet, et chaque partie se réduit à un problème d'algorithmique, d'arithmétique, de graphes,
de probabilités ou de bases de données.

Le sujet comporte \textbf{cinq parties largement indépendantes}, de difficulté croissante.
Les fonctions écrites dans une question peuvent être librement réutilisées ensuite.

\subsection*{Notations et conventions}

\begin{itemize}[leftmargin=1.4em,itemsep=2pt]
\item Tous les entiers manipulés sont des entiers Python (précision arbitraire). Les
opérateurs binaires \code{&} (et), \code{|} (ou), \code{^} (ou exclusif), \code{<<} et
\code{>>} (décalages) sont autorisés et supposés en $O(1)$ sur les entiers du sujet.
\item Une \textbf{adresse IPv4} est un entier de $32$ bits, noté aussi sous forme
\emph{décimale pointée} $\texttt{a.b.c.d}$ où $a,b,c,d\in[\![0,255]\!]$, représentant
l'entier $a\cdot 2^{24}+b\cdot 2^{16}+c\cdot 2^{8}+d$.
\item On note $[\![m,n]\!]$ l'ensemble des entiers $k$ tels que $m\le k\le n$.
\item En Python, seules les constructions du programme officiel sont autorisées~: types de
base, listes, tuples, dictionnaires, ensembles, boucles, fonctions, récursivité, classes
simples, ainsi que les modules \code{math}, \code{random} et \code{heapq}. Les fonctions
\code{sorted} et \code{sort} sont autorisées~; on admettra qu'elles trient $m$ éléments en
$O(m\log m)$.
\item On admet que les opérations élémentaires sur un dictionnaire ou un ensemble
(lecture, écriture, test d'appartenance) s'effectuent en $O(1)$ \emph{en moyenne}.
\item Sauf mention contraire, les complexités demandées sont des complexités \emph{en
temps} dans le pire des cas.
\end{itemize}

\bigskip
\noindent\emph{Le sujet comporte 5 parties et 60 questions. Les parties sont largement
indépendantes~; à l'intérieur d'une partie, les questions s'enchaînent.}

\newpage

%=======================================================================
\section*{Partie I --- Adressage IPv4 et agrégation CIDR}
%=======================================================================

\emph{Références~: \rfc{791} (\emph{Internet Protocol}, 1981), \rfc{4632} (\emph{Classless
Inter-domain Routing}, 2006).}

\medskip
Un \textbf{préfixe} CIDR est un couple $(p,n)$, noté $p/n$, où $n\in[\![0,32]\!]$ est la
\emph{longueur} du préfixe et $p$ une adresse dont les $32-n$ bits de poids faible sont
nuls. Il désigne l'ensemble
\[ \mathcal{A}(p/n) \;=\; \{\,a\in[\![0,2^{32}-1]\!] \ :\ \text{$a$ et $p$ ont les mêmes
$n$ bits de poids fort}\,\}. \]
Par exemple $\texttt{192.168.1.0/24}$ désigne les $256$ adresses de
$\texttt{192.168.1.0}$ à $\texttt{192.168.1.255}$.

\Q Écrire deux fonctions \code{ip_vers_entier(s)} et \code{entier_vers_ip(x)} qui convertissent
une chaîne en décimale pointée en l'entier correspondant, et réciproquement. Préciser leur
complexité.

\Q Écrire une fonction \code{masque(n)} renvoyant, \textbf{sans boucle}, l'entier dont les
$n$ bits de poids fort valent $1$ et les autres $0$. Démontrer que
$\texttt{masque}(n)=2^{32}-2^{32-n}$.

\Q Écrire \code{adresse_reseau(a, n)} et \code{diffusion(a, n)} renvoyant respectivement la
plus petite et la plus grande adresse du préfixe de longueur $n$ contenant $a$. Démontrer
que $\texttt{diffusion}(a,n)=\texttt{adresse\_reseau}(a,n)+2^{32-n}-1$.

\Q Dans un réseau IPv4, l'adresse de réseau et l'adresse de diffusion ne sont pas
attribuables à une machine. En déduire le nombre d'adresses utilisables dans un préfixe de
longueur $n\le 30$, et démontrer la formule. Que se passe-t-il pour $n=31$ et $n=32$~?
Commenter.

\Q Écrire \code{dans_prefixe(a, p, n)} testant si $a\in\mathcal{A}(p/n)$, en utilisant
\textbf{un seul} décalage et \textbf{un seul} ou exclusif. Démontrer l'équivalence
\[ a\in\mathcal{A}(p/n) \iff (a \oplus p)\ \texttt{>>}\ (32-n) = 0 \qquad (n\ge 1), \]
où $\oplus$ désigne le ou exclusif bit à bit.

\Q \textbf{(Agrégation.)} Soient $p/n$ et $q/n$ deux préfixes distincts de même longueur
$n\ge 1$. Démontrer qu'il existe un préfixe $r/(n-1)$ tel que
$\mathcal{A}(r/(n-1))=\mathcal{A}(p/n)\cup\mathcal{A}(q/n)$ \emph{si et seulement si} $p$ et
$q$ ne diffèrent que par le bit de rang $32-n$. Expliquer en une phrase l'intérêt de ce
résultat pour la taille des tables de routage.

\Q \textbf{(Famille laminaire.)} Démontrer que deux préfixes CIDR quelconques $p/n$ et $q/m$
vérifient exactement l'une des trois propriétés suivantes~: $\mathcal{A}(p/n)\subseteq
\mathcal{A}(q/m)$, ou $\mathcal{A}(q/m)\subseteq\mathcal{A}(p/n)$, ou
$\mathcal{A}(p/n)\cap\mathcal{A}(q/m)=\varnothing$. \emph{(Ce résultat est utilisé en
partie~III.)}

\Q Applications numériques. Déterminer, en justifiant~:
\begin{enumerate}[label=(\alph*),leftmargin=2em,itemsep=1pt]
\item le masque de $\texttt{/12}$ en décimale pointée~;
\item l'adresse de réseau de $\texttt{172.16.35.7/12}$~;
\item l'adresse de réseau, l'adresse de diffusion et le nombre d'adresses utilisables de
$\texttt{192.168.1.130/26}$.
\end{enumerate}

%=======================================================================
\section*{Partie II --- La somme de contrôle Internet}
%=======================================================================

\emph{Références~: \rfc{1071} (\emph{Computing the Internet Checksum}, 1988), \rfc{1624}
(\emph{Computation of the Internet Checksum via Incremental Update}, 1994).}

\medskip
Les protocoles IP, TCP et UDP protègent leurs en-têtes par une \textbf{somme de contrôle}
de $16$ bits reposant sur l'\emph{addition en complément à un}. On note
$N=2^{16}=65536$ et l'on travaille sur $[\![0,N-1]\!]$.

\begin{defi}
L'\textbf{addition en complément à un} de $a,b\in[\![0,N-1]\!]$, notée $a\oplus' b$, est
obtenue en additionnant $a$ et $b$ puis en \emph{réinjectant la retenue} dans les bits de
poids faible (\emph{end-around carry})~:
\[ a\oplus' b \;=\; \big((a+b)\bmod N\big) \;+\; \big\lfloor (a+b)/N \big\rfloor . \]
Le \textbf{complément} de $x$ est $\overline{x}=x\oplus \texttt{0xFFFF}$ (ou exclusif), soit
$N-1-x$.
\end{defi}

\Q Écrire une fonction \code{add16(a, b)} calculant $a\oplus' b$ en une ligne. Démontrer que
le résultat appartient bien à $[\![0,N-1]\!]$ (on prendra garde au cas où la réinjection
provoque à son tour une retenue).

\Q Démontrer que pour tous $a,b\in[\![0,N-1]\!]$ on a $a\oplus' b\equiv a+b \pmod{N-1}$.
En déduire que $\oplus'$ est \textbf{commutative} et \textbf{associative}, et que $0$ en est
un élément neutre. Que vaut $x\oplus'\overline{x}$~? Commenter l'existence de \emph{deux}
représentations de zéro, $\texttt{0x0000}$ et $\texttt{0xFFFF}$.

\Q Le calcul de la somme de contrôle d'une suite d'octets se fait en trois temps, que l'on
demande de programmer séparément.

\begin{enumerate}[label=(\alph*),leftmargin=2em,itemsep=4pt]
\item \textbf{Découpage en mots.} Les octets sont regroupés \emph{deux par deux dans l'ordre
de lecture}, le premier octet de chaque paire étant celui de \emph{poids fort}~: la paire
d'octets $(u,v)$ donne le mot $256\,u+v$. Si le nombre d'octets est impair, on complète par
un octet nul \emph{à droite} (ce bourrage sert uniquement au calcul et ne fait pas partie du
paquet transmis).

Écrire \code{mots16(octets)} renvoyant la liste des mots obtenus. Vérifier que
\code{mots16([0x45, 0x00, 0x00, 0x73])} vaut \code{[0x4500, 0x0073]}.

\item \textbf{Somme.} Écrire \code{somme_un(mots)} renvoyant l'addition en complément à un de
tous les mots de la liste, effectuée de gauche à droite en partant de $0$~:
\[ \texttt{somme\_un}([m_0,\dots,m_{\ell-1}]) \;=\;
   \big(\cdots\big((0\oplus' m_0)\oplus' m_1\big)\cdots\big)\oplus' m_{\ell-1}. \]
La question 10 assure que le résultat ne dépend ni du parenthésage ni de l'ordre des mots~;
on ne demande pas de le rejustifier.

\item \textbf{Complément.} Écrire \code{checksum(octets)} renvoyant le complément de la somme
précédente, c'est-à-dire
$\overline{\texttt{somme\_un}(\texttt{mots16}(\texttt{octets}))}$. Vérifier sur la liste de
quatre octets ci-dessus que l'on obtient \code{0xBA8C}.
\end{enumerate}

Donner, pour chacune de ces trois fonctions, la complexité en temps \emph{et} en espace en
fonction du nombre $k$ d'octets. Indiquer enfin comment fusionner \code{mots16} et
\code{somme_un} en une seule boucle afin de ramener l'espace auxiliaire à $O(1)$, et écrire
la fonction correspondante.

\Q Le champ « somme de contrôle » est mis à zéro pour le calcul, puis remplacé par la valeur
obtenue. Démontrer que le récepteur, en calculant \code{somme_un} sur le paquet
\emph{complet} (champ de contrôle inclus), doit trouver $\texttt{0xFFFF}$. Écrire la fonction
\code{verifie(octets)} correspondante.

\Q \textbf{(Indépendance de l'ordre des octets, \rfc{1071} §2.B.)} On note
$\sigma(x)$ le mot obtenu en échangeant les deux octets de $x$. Démontrer que
\[ \sigma(x_1)\oplus'\cdots\oplus'\sigma(x_k) \;=\; \sigma\big(x_1\oplus'\cdots\oplus'
x_k\big). \]
Quelle conséquence pratique cela a-t-il pour une machine dont l'ordre des octets natif
diffère de celui du réseau~?

\Q Démontrer qu'une erreur de transmission portant sur \textbf{un seul bit} est toujours
détectée.

\Q Montrer en revanche que la permutation de deux mots de $16$ bits n'est \emph{jamais}
détectée, et exhiber une erreur portant sur \emph{deux} bits qui passe inaperçue. Quelle
propriété de $\oplus'$ est en cause~? Comparer brièvement avec un code polynomial de type
CRC.

\Q \textbf{(Mise à jour incrémentale.)} Un routeur qui décrémente le champ « durée de vie »
d'un paquet n'a pas besoin de recalculer toute la somme. On note $HC$ l'ancienne somme de
contrôle, $m$ l'ancienne valeur du mot modifié, $m'$ la nouvelle, et $HC'$ la nouvelle somme.
\begin{enumerate}[label=(\alph*),leftmargin=2em,itemsep=1pt]
\item La \rfc{1141} proposait $HC' = HC\oplus' m\oplus'\overline{m'}$. Implémenter cette
formule dans une fonction \code{maj_1141(HC, m, mp)}.
\item La \rfc{1624} la corrige en $HC' = \overline{\ \overline{HC}\oplus'\overline{m}\oplus'
m'\ }$. Implémenter \code{maj_1624(HC, m, mp)}.
\item On suppose que la somme en complément à un de tous les \emph{autres} mots de l'en-tête
vaut $C_{\text{autres}}=\texttt{0xCD7A}$, que $m=\texttt{0x5555}$ et
$m'=\texttt{0x3285}$. Calculer $HC$, puis $HC'$ par recalcul complet, puis par chacune des
deux formules. Conclure.
\item Expliquer précisément \emph{pourquoi} la formule de la \rfc{1141} échoue. \emph{On
pourra invoquer la réponse à la question 10.}
\end{enumerate}

\Q Application numérique. On considère l'en-tête IPv4 suivant, donné en hexadécimal, dont le
champ de contrôle (octets d'indices $10$ et $11$) a été mis à zéro~:
\begin{center}\ttfamily\small
45 00 00 73 00 00 40 00 40 11 00 00 c0 a8 00 01 c0 a8 00 c7
\end{center}
Calculer la somme de contrôle à y inscrire. Vérifier ensuite que le paquet complet satisfait
bien le test de la question 12.

%=======================================================================
\section*{Partie III --- Routage : plus long préfixe et plus courts chemins}
%=======================================================================

\emph{Références~: \rfc{4632} (CIDR), \rfc{2328} (OSPF, algorithme à état de liens),
\rfc{2453} (RIP, algorithme à vecteur de distance).}

\subsection*{III.1 --- Recherche du plus long préfixe correspondant}

Une \textbf{table de routage} est une liste de triplets $(p,n,\texttt{saut})$. Pour router
une adresse $a$, le routeur retient le préfixe \emph{le plus long} parmi ceux qui contiennent
$a$, et transmet le paquet au saut associé. Le préfixe $\texttt{0.0.0.0/0}$, présent dans
toute table, joue le rôle de route par défaut.

\Q Écrire \code{plus_long_prefixe(table, a)} par parcours exhaustif de la table. Donner sa
complexité en fonction du nombre $k$ d'entrées.

\Q Démontrer, à l'aide de la question 7, que le préfixe le plus long contenant $a$ est
\textbf{unique} dès lors que la table ne contient pas deux entrées de même préfixe.

\Q On représente maintenant la table par un \textbf{arbre binaire de préfixes} (ou
\emph{trie})~: un n\oe ud à la profondeur $d$ représente une suite de $d$ bits~; son fils
gauche (resp. droit) représente cette suite suivie du bit $0$ (resp. $1$). Un n\oe ud porte
un saut si et seulement si la suite de bits correspondante est un préfixe de la table.
Écrire une méthode \code{inserer(p, n, saut)}. Donner la complexité en temps d'une insertion
et la complexité en espace de la structure pour $k$ entrées.

\Q Écrire la méthode \code{chercher(a)} qui descend dans l'arbre en suivant les bits de $a$
et renvoie le saut du \emph{dernier} n\oe ud rencontré qui en porte un. Énoncer l'invariant
de la boucle et démontrer la correction.

\Q Donner la complexité de \code{chercher}. Comparer avec la question 18 et commenter la
pertinence de cette structure pour un routeur de c\oe ur de réseau, dont la table comporte
aujourd'hui environ $10^6$ entrées.

\subsection*{III.2 --- Plus courts chemins dans le graphe du réseau}

On modélise le réseau par un graphe non orienté pondéré $G=(S,A)$~: les sommets sont les
routeurs, numérotés de $0$ à $V-1$, et le poids d'une arête est le \emph{coût} du lien
(entier strictement positif). On note $E=|A|$.

\Q Écrire \code{dijkstra(V, adj, s)} calculant les distances de $s$ à tous les sommets, où
\code{adj[u]} est la liste des couples $(v,w)$. On utilisera le module \code{heapq}. Donner
la complexité.

\Q Démontrer la correction de l'algorithme de Dijkstra en énonçant l'invariant vérifié à
chaque extraction du tas. Où l'hypothèse « poids positifs » est-elle utilisée~?

\Q Écrire \code{bellman_ford(V, aretes, s)}, où \code{aretes} est la liste des arcs
$(u,v,w)$. La fonction renverra \code{None} si le graphe contient un circuit de poids
strictement négatif. Donner la complexité.

\Q Démontrer, par récurrence sur $i$, qu'après $i$ tours de la boucle externe,
$\texttt{d[v]}$ est majoré par le poids du plus court chemin de $s$ à $v$ \emph{utilisant au
plus $i$ arêtes}. En déduire la correction et justifier le nombre $V-1$ de tours.

\Q \textbf{(Comptage à l'infini.)} Le protocole RIP est un algorithme à vecteur de distance~:
chaque routeur ne connaît que les distances annoncées par ses voisins. Sur le graphe réduit
à trois routeurs $A - B - C$ en ligne (arêtes de poids $1$), décrire ce qui se produit
lorsque le lien $B-C$ tombe, et expliquer pourquoi les distances annoncées croissent
indéfiniment. Quelle contre-mesure élémentaire RIP adopte-t-il~? Commenter son coût.

\Q Application numérique. On considère le graphe non orienté à $6$ sommets d'arêtes
\[ (0,1,7)\ (0,2,9)\ (0,5,14)\ (1,2,10)\ (1,3,15)\ (2,3,11)\ (2,5,2)\ (3,4,6)\ (4,5,9). \]
Calculer, en détaillant les étapes de Dijkstra, les distances depuis le sommet $0$, ainsi
qu'un plus court chemin de $0$ à $4$.

%=======================================================================
\section*{Partie IV --- TCP : fiabilité, temporisateur et congestion}
%=======================================================================

\emph{Références~: \rfc{9293} (TCP), \rfc{6298} (\emph{Computing TCP's Retransmission
Timer}, 2011), \rfc{5681} (\emph{TCP Congestion Control}, 2009), V.~Jacobson,
\emph{Congestion Avoidance and Control}, SIGCOMM~1988~; M.~Mathis, J.~Semke, J.~Mahdavi,
T.~Ott, \emph{The Macroscopic Behavior of the TCP Congestion Avoidance Algorithm},
ACM~CCR~27(3), 1997.}

\subsection*{IV.1 --- Numéros de séquence circulaires}

TCP numérote les octets par un compteur de $32$ bits qui \emph{boucle}~: on travaille dans
$\mathbb{Z}/M\mathbb{Z}$ avec $M=2^{32}$. On veut décider si un numéro $a$ « précède » un
numéro $b$.

\Q On pose $\texttt{avant}(a,b) = \big[(b-a)\bmod M < 2^{31}\ \text{et}\ a\ne b\big]$.
Écrire la fonction. Démontrer que si tous les numéros en circulation à un instant donné
appartiennent à une fenêtre de largeur $W$, alors \code{avant} rend le verdict correct dès
que $W<2^{31}$, et exhiber une ambiguïté lorsque $W=2^{31}$.

\subsection*{IV.2 --- Réassemblage des segments}

Un segment reçu couvre l'intervalle d'octets $[\,d,f\,[$. Les segments arrivent dans le
désordre et peuvent se chevaucher. Le récepteur doit maintenir la liste des zones
effectivement reçues.

\Q Écrire \code{fusionner(segments)} qui, à partir d'une liste de couples $(d,f)$, renvoie la
liste des intervalles maximaux couverts, triée par début croissant. Donner la complexité en
fonction du nombre $k$ de segments.

\Q Démontrer la correction de \code{fusionner}~: la réunion des intervalles renvoyés est
égale à celle des intervalles fournis, et les intervalles renvoyés sont deux à deux disjoints
et \emph{non adjacents}. On énoncera un invariant de boucle.

\subsection*{IV.3 --- Estimation du temporisateur de retransmission}

Le \rfc{6298} prescrit l'algorithme suivant, dû à Jacobson et Karels. À la première mesure
de temps d'aller-retour $R$~:
\[ \mathrm{SRTT}\leftarrow R,\qquad \mathrm{RTTVAR}\leftarrow R/2, \]
puis, à chaque mesure ultérieure $R'$, \emph{dans cet ordre}~:
\[ \mathrm{RTTVAR}\leftarrow(1-\beta)\,\mathrm{RTTVAR}+\beta\,|\mathrm{SRTT}-R'|,\qquad
   \mathrm{SRTT}\leftarrow(1-\alpha)\,\mathrm{SRTT}+\alpha\,R', \]
avec $\alpha=1/8$ et $\beta=1/4$~; enfin
$\mathrm{RTO}\leftarrow\max\big(1,\ \mathrm{SRTT}+4\,\mathrm{RTTVAR}\big)$ (en secondes).

\Q Écrire deux fonctions \code{rto_initial(R)} et \code{rto_suivant(SRTT, RTTVAR, R)}
implémentant ces règles. Pourquoi l'ordre des deux affectations est-il important~?

\Q On note $R_0$ la première mesure et $R_1,\dots,R_n$ les suivantes, et $S_n$ la valeur de
$\mathrm{SRTT}$ après la $n$-ième mise à jour. Démontrer par récurrence que
\[ S_n \;=\; (1-\alpha)^n R_0 \;+\; \alpha\sum_{i=1}^{n}(1-\alpha)^{\,n-i}R_i . \]
Vérifier que la somme des coefficients vaut exactement $1$.

\Q On appelle \emph{mémoire effective} de l'estimateur le nombre $\mu=\sum_{j\ge 0}(j+1)\,
\alpha(1-\alpha)^j$ d'échantillons « équivalents » qu'il retient. Calculer $\mu$ en fonction
de $\alpha$, puis sa valeur pour $\alpha=1/8$. Déterminer également le plus petit entier $j$
tel que le poids $\alpha(1-\alpha)^j$ soit inférieur au dixième du poids le plus récent.
Interpréter.

\Q \textbf{(Algorithme de Karn et repli exponentiel.)} Lorsqu'un segment est retransmis, on
ne peut pas savoir si l'accusé de réception répond à la première ou à la seconde émission.
Expliquer en quoi mesurer $R$ dans ce cas biaiserait l'estimateur, et dans quel sens. Le
\rfc{6298} impose de plus $\mathrm{RTO}\leftarrow 2\,\mathrm{RTO}$ à chaque expiration.
Calculer le temps d'attente cumulé après $n$ expirations successives à partir d'un
$\mathrm{RTO}$ initial $T_0$, et donner sa valeur pour $T_0=1$~s et $n=5$.

\Q On suppose que chaque transmission d'un segment échoue indépendamment avec probabilité
$p\in\,]0,1[$. Soit $N$ le nombre total de transmissions nécessaires. Reconnaître la loi de
$N$, donner $\mathbb{P}(N=k)$, $\mathbb{P}(N>k)$ et $\mathbb{E}[N]$. Démontrer que la
transmission aboutit \textbf{presque sûrement}. Application numérique pour $p=0{,}1$.

\subsection*{IV.4 --- Contrôle de congestion : le débit d'une connexion TCP}

TCP règle son débit par une \emph{fenêtre de congestion} $W$ (en segments). En régime
d'évitement de congestion (\rfc{5681}), $W$ augmente de $1$ à chaque aller-retour~; à chaque
perte, $W$ est divisée par $2$. C'est le schéma dit \textbf{AIMD} (\emph{additive increase,
multiplicative decrease}). On suppose le régime périodique~: $W$ croît de $W_{\max}/2$ à
$W_{\max}$, une perte survient, et le cycle recommence. On note $\mathrm{RTT}$ la durée d'un
aller-retour, supposée constante, et $W_{\max}$ pair.

\Q Démontrer que le nombre de segments émis pendant un cycle vaut exactement
\[ N_{\text{cycle}} \;=\; \frac{3W_{\max}^{2}}{8}-\frac{W_{\max}}{4}, \]
et que la fenêtre moyenne sur un cycle vaut $\overline{W}=\dfrac{3W_{\max}-2}{4}$.

\Q On suppose qu'exactement une perte se produit par cycle, et l'on note $p$ la probabilité
de perte d'un segment, de sorte que $p=1/N_{\text{cycle}}$. En négligeant les termes
d'ordre~$1$ devant les termes d'ordre~$2$, démontrer la \textbf{formule de Mathis}~:
\[ W_{\max}\approx\sqrt{\frac{8}{3p}}\,,\qquad\text{d'où}\qquad
   \text{débit}\;\approx\;\frac{\mathrm{MSS}}{\mathrm{RTT}}\sqrt{\frac{3}{2p}}
   \;\approx\;\frac{1{,}22\ \mathrm{MSS}}{\mathrm{RTT}\sqrt{p}}\quad\text{(octets/s)}, \]
où $\mathrm{MSS}$ est la taille d'un segment en octets.

\Q Application numérique~: $\mathrm{MSS}=1460$ octets, $\mathrm{RTT}=100$~ms. Calculer le
débit prévu et la fenêtre $W_{\max}$ pour $p=10^{-2}$ puis pour $p=10^{-4}$.

\Q Commenter la formule de Mathis. En particulier~: (a) quel est l'effet d'une division par
$4$ du taux de perte~? (b) deux connexions partageant le même lien mais de $\mathrm{RTT}$
différents obtiennent-elles la même part du débit~? Quelle conséquence pour l'équité~?

%=======================================================================
\section*{Partie V --- HTTP, cache et HTTPS}
%=======================================================================

\emph{Références~: \rfc{9110} et \rfc{9112} (HTTP/1.1, 2022), \rfc{9111} (cache HTTP),
\rfc{8446} (TLS~1.3, 2018)~; W.~Diffie et M.~Hellman, \emph{New Directions in Cryptography},
IEEE~IT-22(6), 1976.}

\subsection*{V.1 --- Analyse d'une requête HTTP}

Une requête HTTP/1.1 est un texte de la forme
\begin{center}\ttfamily\small
\begin{tabular}{l}
GET /index.html?p=1 HTTP/1.1\textbackslash r\textbackslash n\\
Host: cpgeacademy.org\textbackslash r\textbackslash n\\
Content-Length: 5\textbackslash r\textbackslash n\\
\textbackslash r\textbackslash n\\
HELLO
\end{tabular}
\end{center}
La première ligne (\emph{ligne de requête}) contient trois champs séparés par une espace.
Suivent des \emph{en-têtes} de la forme \code{Nom: valeur}, un par ligne. Une ligne vide,
c'est-à-dire la séquence $\texttt{CRLF}\,\texttt{CRLF}$ (soit
\code{"\\r\\n\\r\\n"}), sépare l'en-tête du corps.

\Q Écrire \code{analyser_requete(texte)} renvoyant le quintuplet
\[ \big(\text{méthode},\ \text{cible},\ \text{version},\
   \text{dictionnaire des en-têtes},\ \text{corps}\big), \]
les noms d'en-têtes étant mis en minuscules. Donner la complexité.

\Q Un serveur lit le flux \emph{caractère par caractère} et doit détecter la fin de l'en-tête
sans jamais revenir en arrière. Construire un \textbf{automate fini déterministe} à quatre
états reconnaissant la première occurrence de $\texttt{CRLFCRLF}$. Donner sa table de
transition, puis écrire la fonction \code{fin_entete(flux)} renvoyant l'indice suivant la fin
de l'en-tête, ou $-1$.

\Q Démontrer la correction de cet automate en énonçant l'invariant~: \emph{à chaque instant,
l'état courant est la longueur du plus long suffixe du texte déjà lu qui soit un préfixe de}
$\texttt{CRLFCRLF}$. En déduire que l'analyse est linéaire en la longueur du flux et ne
nécessite aucun retour en arrière. Vérifier l'invariant sur l'entrée
\code{"\\r\\n\\r\\r\\n\\r\\n"}.

\Q Lorsque la taille du corps n'est pas connue à l'avance, HTTP/1.1 utilise le codage
\emph{par morceaux}~: le corps est une suite de blocs, chacun précédé de sa taille écrite en
hexadécimal et suivi de $\texttt{CRLF}$~; un bloc de taille nulle termine le corps. Écrire
\code{decoder_chunked(s)}. Exhiber un variant de boucle et en déduire la terminaison. Donner
la complexité.

\subsection*{V.2 --- Le cache}

\Q Un cache de capacité $C$ objets applique la politique \textbf{LRU} (\emph{least recently
used})~: en cas de saturation, l'objet dont le dernier accès est le plus ancien est évincé.
Proposer une structure de données permettant de traiter un accès en $O(1)$ \emph{en moyenne},
et écrire la classe correspondante avec une méthode \code{acceder(cle)} renvoyant \code{True}
en cas de succès. Justifier la complexité et énoncer l'invariant reliant la structure à
l'ordre des derniers accès.

\Q Dérouler la politique LRU de capacité $3$ sur la suite d'accès
\[ a,\ b,\ c,\ a,\ d,\ b,\ a,\ c,\ d,\ a \]
en indiquant à chaque étape le contenu du cache, et donner le nombre de succès.

\subsection*{V.3 --- Choix optimal du contenu d'un cache (programmation dynamique)}

Un serveur dispose de $n$ objets~; l'objet $i$ occupe $t_i$ octets et est demandé $f_i$ fois
par heure. On veut choisir un sous-ensemble $S$ d'objets à placer en mémoire, de taille
totale au plus $C$, maximisant le nombre de requêtes servies depuis le cache~:
\[ \mu(t,f,C) \;=\; \max\Big\{\textstyle\sum_{i\in S}f_i \ :\ S\subseteq\{0,\dots,n-1\},\
\sum_{i\in S}t_i\le C\Big\}. \]

\Q On pose, pour $0\le i\le n$ et $0\le j\le C$, $m_i(j)$ le nombre maximal de requêtes
servies par un sous-ensemble de $\{0,\dots,i-1\}$ de taille totale au plus $j$. Donner
$m_0$, puis établir la relation de récurrence exprimant $m_{i+1}$ en fonction de $m_i$.

\Q Écrire \code{cache_optimal(t, f, C)} implémentant cette récurrence avec un \textbf{unique}
tableau de taille $C+1$ mis à jour en place. Justifier précisément le \textbf{sens de
parcours} de $j$, et montrer sur un exemple explicite ce que renverrait l'algorithme si l'on
parcourait dans l'autre sens.

\Q Démontrer la correction en énonçant et en prouvant, par récurrence sur le nombre $i$
d'objets traités, l'\textbf{invariant} vérifié par le tableau à la fin de la $i$-ième
itération de la boucle externe. Justifier la terminaison.

\Q Donner la complexité en temps et en espace. Le problème de décision associé est
\textsc{Sac-à-Dos}, connu NP-complet~; montrer qu'il appartient à NP en explicitant un
certificat et un vérificateur. Expliquer pourquoi l'algorithme précédent ne contredit pas
cette NP-complétude, et ce qui se produit si les tailles sont exprimées en octets et
$C=10$~Go.

\Q Application numérique. On prend $t=[3,4,5,9,4]$, $f=[60,50,70,90,30]$ (tailles en Mo) et
$C=10$. Calculer $\mu$ et donner un sous-ensemble optimal. Détailler le contenu du tableau à
la fin de chaque itération pour la version réduite $t=[3,4,5]$, $f=[60,50,70]$, $C=8$.

\subsection*{V.4 --- Analyse des journaux du serveur (SQL)}

Le serveur enregistre ses requêtes dans une base relationnelle de schéma~:
\begin{center}\ttfamily\small
\begin{tabular}{l}
journal(\underline{id}, ip, horodatage, methode, chemin, statut, octets, duree\_ms)\\
client(\underline{ip}, pays, organisation)
\end{tabular}
\end{center}
où \code{horodatage} est une chaîne au format \code{'AAAA-MM-JJ HH:MM:SS'} et \code{statut}
le code de réponse HTTP ($2xx$~: succès, $3xx$~: redirection ou contenu inchangé, $4xx$~:
erreur du client, $5xx$~: erreur du serveur). On rappelle que \code{substr(s, i, k)} extrait
$k$ caractères de $s$ à partir de la position $i$ (comptée à partir de $1$).

\Q Écrire une requête SQL donnant, pour chaque méthode, le nombre de requêtes et le volume
total d'octets transmis, par nombre de requêtes décroissant.

\Q Écrire une requête donnant les trois chemins les plus demandés, avec leur nombre de
requêtes.

\Q Écrire une requête donnant, pour chaque heure de la journée, le nombre total de requêtes,
le nombre de réponses $5xx$ et le taux d'erreur en pourcentage arrondi au dixième.

\Q Écrire une requête donnant les adresses IP ayant provoqué au moins trois réponses $4xx$,
avec le nombre correspondant.

\Q Écrire une requête donnant, par pays, le nombre de requêtes et le volume d'octets. Écrire
enfin une requête donnant la liste des chemins qui n'ont \emph{jamais} été servis avec le
statut $200$.

\subsection*{V.5 --- HTTPS : l'établissement du secret partagé}

TLS~1.3 (\rfc{8446}) établit un secret partagé par un échange de Diffie--Hellman. On se
donne un nombre premier $p$ et un entier $g\in[\![2,p-2]\!]$, tous deux publics.

\Q Écrire une fonction \code{puissance_mod(a, e, n)} calculant $a^{e}\bmod n$ par
exponentiation rapide, \textbf{sans} utiliser \code{pow(a, e, n)}. Exhiber un variant de
boucle et en déduire la terminaison. Démontrer la correction en énonçant l'invariant
$r\cdot a^{e}\equiv a_0^{e_0}\pmod n$. Démontrer que le nombre de multiplications modulaires
est au plus $2\lfloor\log_2 e\rfloor+2$.

\Q Le client tire $x$ au hasard et envoie $A=g^{x}\bmod p$~; le serveur tire $y$ et envoie
$B=g^{y}\bmod p$. Chacun calcule alors une clé. Écrire les deux calculs et démontrer qu'ils
donnent le même résultat. Application numérique~: $p=23$, $g=5$, $x=6$, $y=15$~; calculer
$A$, $B$ et la clé partagée.

\Q Un attaquant qui observe le réseau connaît $p$, $g$, $A$ et $B$. Le retrouver $x$ est le
problème du \textbf{logarithme discret}. Décrire un algorithme naïf et donner sa complexité.
Décrire ensuite le principe de l'algorithme « pas de bébé, pas de géant » en $O(\sqrt{p})$
opérations et $O(\sqrt{p})$ en espace. Pour $p$ de $2048$ bits, comparer les deux coûts au
nombre d'opérations réalisables en un an par une machine effectuant $10^{9}$ opérations par
seconde. Conclure.

\Q \textbf{Question de synthèse (à traiter en quelques lignes).} TLS~1.3 impose que $x$ et
$y$ soient \emph{éphémères}, c'est-à-dire tirés à nouveau pour chaque connexion et effacés
ensuite, et a supprimé la possibilité d'un échange de clés reposant sur le chiffrement par la
clé publique du serveur. Expliquer ce que l'on gagne ainsi si la clé privée du serveur est
compromise \emph{après} l'interception et l'enregistrement d'une session. En quoi la partie~II
(intégrité) et cette partie~V (confidentialité) répondent-elles à des menaces de nature
différente~?

\vfill
\begin{center}\small\textsc{Fin de l'épreuve}\end{center}
\end{document}
